%==============================================================================
% Chapitre 8 - Pression atmosphérique
%==============================================================================

\section{Introduction}

La pression atmosphérique est l'une des principales caractéristiques de
l'atmosphère terrestre.

Elle influence directement :

\begin{itemize}
\item la densité de l'air ;
\item la vitesse du son ;
\item la traînée aérodynamique ;
\item les performances balistiques.
\end{itemize}

Bien que ses variations soient invisibles à l'œil nu, elles peuvent modifier
de manière mesurable la trajectoire d'un projectile.

\begin{aretenir}

La pression atmosphérique influence principalement la balistique par son effet
sur la densité de l'air.

\end{aretenir}

\section{Qu'est-ce que la pression atmosphérique ?}

L'air possède une masse.

Sous l'effet de la gravité, les couches supérieures de l'atmosphère exercent
une force sur les couches inférieures.

Cette force se traduit par une pression.

La pression atmosphérique correspond donc au poids de la colonne d'air située
au-dessus d'un point donné.

\section{Unités de pression}

Dans le système international, la pression s'exprime en pascals (Pa).

En météorologie, on utilise fréquemment l'hectopascal (hPa).

La relation est :

\begin{equation}
1~\mathrm{hPa}=100~\mathrm{Pa}
\end{equation}

La pression atmosphérique standard au niveau de la mer vaut :

\begin{equation}
P_0 = 101,325~\mathrm{Pa}
\end{equation}

soit :

\begin{equation}
P_0 = 1013,25~\mathrm{hPa}
\end{equation}

\section{Pourquoi la pression existe-t-elle ?}

Imaginons une colonne d'air de section unitaire s'étendant du sol jusqu'au
sommet de l'atmosphère.

Cette colonne possède une masse.

Sous l'effet de la gravité, cette masse exerce une force sur le sol.

La pression atmosphérique est simplement le rapport entre cette force et la
surface considérée.

\section{Variation avec l'altitude}

Lorsque l'altitude augmente, la quantité d'air située au-dessus de
l'observateur diminue.

La pression atmosphérique diminue donc avec l'altitude.

Cette diminution n'est pas linéaire.

Elle est relativement rapide dans les basses couches de l'atmosphère.

\begin{exemple}

Valeurs typiques :

\begin{center}
\begin{tabular}{lr}
\toprule
Altitude & Pression approximative \\
\midrule
0 m & 1013 hPa \\
1000 m & 899 hPa \\
2000 m & 795 hPa \\
3000 m & 701 hPa \\
\bottomrule
\end{tabular}
\end{center}

\end{exemple}

\section{Influence sur la densité}

La pression et la densité sont étroitement liées.

À température constante :

\begin{itemize}
\item une pression plus élevée conduit généralement à une densité plus
élevée ;
\item une pression plus faible conduit généralement à une densité plus
faible.
\end{itemize}

Cette relation explique pourquoi les projectiles rencontrent moins de
résistance à haute altitude.

\section{Influence sur les trajectoires}

Lorsque la pression diminue :

\begin{itemize}
\item la densité de l'air diminue ;
\item la traînée diminue ;
\item le projectile conserve mieux sa vitesse ;
\item le temps de vol peut être réduit.
\end{itemize}

Ces effets deviennent particulièrement visibles aux longues distances.

\begin{exemple}

Deux tirs identiques réalisés :

\begin{itemize}
\item au niveau de la mer ;
\item à 2500 m d'altitude ;
\end{itemize}

produiront généralement des trajectoires différentes.

Le projectile tiré en altitude conservera davantage d'énergie et de vitesse.

\end{exemple}

\section{Systèmes météorologiques}

La pression atmosphérique varie également avec les conditions
météorologiques.

Les systèmes dépressionnaires et anticycloniques modifient localement la
pression.

Ces variations sont généralement plus faibles que celles dues à l'altitude,
mais elles peuvent être mesurables dans les applications balistiques de
précision.

\section{Mesure de la pression}

La pression atmosphérique est mesurée à l'aide d'un baromètre.

Les stations météorologiques modernes fournissent généralement :

\begin{itemize}
\item la pression absolue ;
\item la pression corrigée au niveau de la mer ;
\item l'évolution de la pression dans le temps.
\end{itemize}

Lors d'essais balistiques, il est important de préciser quelle valeur a été
utilisée.

\section{Application aux essais}

La pression atmosphérique fait partie des paramètres qui doivent être
enregistrés lors d'une campagne d'essais.

Cette information permet :

\begin{itemize}
\item de reproduire les conditions de tir ;
\item de comparer plusieurs campagnes ;
\item d'améliorer les corrélations entre essais et simulation.
\end{itemize}

\begin{simulation}

Les modèles balistiques avancés utilisent la pression atmosphérique afin
d'estimer la densité de l'air.

Cette information est ensuite utilisée pour calculer la traînée et les autres
effets aérodynamiques.

\end{simulation}

\begin{rigueur}

Dans un gaz parfait, la pression, la densité et la température sont liées.

Cette relation sera approfondie dans les chapitres suivants consacrés à la
densité de l'air et aux modèles atmosphériques.

\end{rigueur}

\section{À retenir}

\begin{aretenir}

\begin{itemize}
\item La pression atmosphérique est due au poids de l'air.
\item Elle diminue lorsque l'altitude augmente.
\item Elle influence directement la densité de l'air.
\item Une pression plus faible conduit généralement à une traînée plus
faible.
\item La pression doit être relevée lors des essais balistiques.
\end{itemize}

\end{aretenir}

